\section*{Task 0: Rank-pattern setup}

\paragraph{Goal.}
Define the cyclic rank-pattern coordinates in the direct weight-space algebra and
remove the indexing ambiguity coming from the cycle.

\paragraph{Lifted admissible intervals.}
For the weight-space algebra
\[
A_L^{\mathrm{wt}}=kQ/I_L^{\mathrm{wt}},
\]
it is convenient to work on the universal cover of the oriented cycle. Define
the lifted admissible index set
\[
\widetilde{\mathcal I}_L^{\mathrm{wt}}
:=
\{(a,b)\in \mathbb Z^2 \mid 0\le a\le L,\ a\le b\le a+L\}
\setminus
\{(a,a+L)\mid 1\le a\le L\}.
\]
Equivalently, an admissible lifted interval starts at a vertex
\[
a\in\{0,\dots,L\},
\]
has length at most \(L\), and the maximal length \(L\) is allowed only when
\[
a=0.
\]

\paragraph{Dictionary with the cyclic notation.}
Every admissible pair in the current draft can be rewritten uniquely in lifted
coordinates:
\[
(i,j)\in \mathcal I_L^{\mathrm{wt}}
\quad\longleftrightarrow\quad
\begin{cases}
(i,j), & 0\le i\le j\le L,\\
(i,j+L+1), & 1\le i\le L,\ 0\le j\le i-2.
\end{cases}
\]
Thus the lifted interval \([a,b]\) records the unique oriented segment on the
cycle going from \(a\) to \(b \bmod (L+1)\) without repeating vertices.

\paragraph{Lifted path coordinates.}
For \((a,b)\in \widetilde{\mathcal I}_L^{\mathrm{wt}}\), let
\[
q_{ab}
\]
denote the unique nonzero oriented path of \(A_L^{\mathrm{wt}}\) represented by
the interval \([a,b]\) on the universal cover.

Two special cases will be used repeatedly:
\[
q_{0L}=a_L\cdots a_1,
\qquad
q_{L,L+1}=a_0.
\]
The first is the open-chain path corresponding to the end-to-end product
\(\mu(W)\); the second is the closing arrow.

\paragraph{Why this indexing is correct.}
The only oriented paths of length \(L\) in the cycle are
\[
[0,L],\ [1,L+1],\ \dots,\ [L,2L].
\]
The relation ideal \(I_L^{\mathrm{wt}}\) kills every length-\(L\) path passing
through \(a_0\), i.e. every interval \([a,a+L]\) with \(1\le a\le L\), and it
leaves the open-chain path \([0,L]\) untouched. Therefore the nonzero oriented
paths of \(A_L^{\mathrm{wt}}\) are indexed exactly by
\(\widetilde{\mathcal I}_L^{\mathrm{wt}}\).

\paragraph{Lifted rank pattern.}
Let \(M\) be a representation of \(A_L^{\mathrm{wt}}\) with dimension vector
\[
d=(d_0,\dots,d_L).
\]
For each \((a,b)\in \widetilde{\mathcal I}_L^{\mathrm{wt}}\), define
\[
\rho_{ab}(M):=\operatorname{rank}\bigl(q_{ab}(M)\bigr).
\]
The collection
\[
\rho(M):=(\rho_{ab}(M))_{(a,b)\in \widetilde{\mathcal I}_L^{\mathrm{wt}}}
\]
is the lifted cyclic rank pattern of \(M\).

\paragraph{Interpretation in DLN language.}
If \(W\in \Gamma_{d,S}\) and \(\phi(W)\) is the corresponding fixed-arrow
representation, then
\[
\rho_{0L}(W)=\operatorname{rank}(W_L\cdots W_1)=\operatorname{rank}(\mu(W)),
\]
and
\[
\rho_{L,L+1}(W)=\operatorname{rank}(W_0).
\]
These are the two scalar coordinates that encode the end-to-end rank and the
fixed closing-arrow rank.

\paragraph{Status.}
Completed. The lifted indexing resolves the cyclic bookkeeping cleanly and will
be used in every subsequent task of Conjecture 3.
